\Aufgabe{Neues Projekt mit BlueJ-Blockly}

1. BlueJ öffnen, neues Projekt erstellen, mit 'Choose' einen Speicherort auswählen und Namen eingeben. Immer auf Laufwerk H:/ speichern!
\\\\

\includegraphics[width=0.15\textwidth]{Aufgaben/BlocksAndJava/img/0_bluej.png} \hfill
\includegraphics[width=0.4\textwidth]{Aufgaben/BlocksAndJava/img/1_newProject.png} \hfill
\includegraphics[width=0.4\textwidth]{Aufgaben/BlocksAndJava/img/2_projectName.png}

\vspace{0.4cm}
2. Neue Klasse erstellen und per Rechtsklick mit Blockly öffnen.

\includegraphics[width=0.4\textwidth]{Aufgaben/BlocksAndJava/img/3_newClass.png} \hspace{0.5cm}
\includegraphics[width=0.25\textwidth]{Aufgaben/BlocksAndJava/img/4_blockly.png}

\vspace{0.4cm}
3. In Blockly (Fenster falls nötig in Taskleiste anklicken) neuen Methodenblock einfügen und falls benötigt Parameter (= input) über das Zahnrad hinzufügen. Parameter könne über Variablen verwendet werden.

\includegraphics[width=0.3\textwidth]{Aufgaben/BlocksAndJava/img/5_createParams.png} \hspace{0.5cm}
\includegraphics[width=0.3\textwidth]{Aufgaben/BlocksAndJava/img/6_variables.png}

\vspace{0.4cm}
4. In BlueJ die Klasse übersetzen/compilen, neues Objekt (mit beliebigem Namen) erzeugen und mit Rechtsklick auf das Objekt Methode aufrufen.

\centering
\includegraphics[width=0.25\textwidth]{Aufgaben/BlocksAndJava/img/7_compile.png} \hfill
\includegraphics[width=0.25\textwidth]{Aufgaben/BlocksAndJava/img/8_newObject.png} \hfill
\includegraphics[width=0.4\textwidth]{Aufgaben/BlocksAndJava/img/9_runMethod.png}